\documentclass{article}
\usepackage[T1]{fontenc}
\usepackage{lmodern}
\usepackage{graphicx}
\usepackage{amsmath,amssymb}
\usepackage[a4paper,margin=2cm]{geometry}
\usepackage[absolute,overlay]{textpos}
\setlength{\TPHorizModule}{1mm}
\setlength{\TPVertModule}{1mm}
\usepackage[indonesian]{babel}
\usepackage{hyperref}
\setlength{\emergencystretch}{2em}
% unit: Halaman 17

\begin{document}

% === Halaman 17 (python/native) ===

Contoh 4:  Misalkan 



Plainteks: 011000110101101110


 dan kunci publik adalah hasil dari Contoh 3,



Kunci publik = \{62, 93, 81, 88, 102, 37\},  


Kunci privat adalah \{2, 3, 6, 13, 27, 52\}.

 Plainteks dibagi menjadi blok yang panjangnya 6, kemudian setiap 

bit di dalam blok dikalikan dengan elemen yang berkorepsonden di 

dalam kunci publik:

17

\end{document}
